% main.tex — объединенный файл документа
\documentclass[12pt,a4paper]{article}

% =============================================================================
% КОНФИГУРАЦИЯ И ПРЕАМБУЛА
% =============================================================================

% --- конфигурация метаданных ---
% config.tex — метаданные и макросы (прямые определения)
\ProvidesFile{config.tex}[2025/10/24 MUFSIPA config]

% ---------- Прямые значения (редактируйте по необходимости) ----------
\newcommand{\Company}{Компания}
\newcommand{\DocumentTitleMain}{Шаблончик}
\newcommand{\DocumentTitleSub}{для оформления документиков}
\newcommand{\Version}{1.0}
\newcommand{\DocDate}{17.10.2025}
\newcommand{\AuthorName}{Фамилия И.О.}
\newcommand{\AgreedBy}{Фамилия И.О.}
\newcommand{\ApprovedBy}{Фамилия И.О.}  % если пусто, оставьте как есть
\newcommand{\City}{Москва}

% --- преамбула и стили ---
% preamble.tex — пакеты и глобальные установки, соблюдающие стиль документа
\usepackage[utf8]{inputenc}
\usepackage[T2A]{fontenc}
\usepackage[russian]{babel}

% Times New Roman-like
\usepackage{newtxtext,newtxmath} % похож на Times; обеспечивает текст+математику
\usepackage{setspace}
\onehalfspacing % межстрочный интервал 1.5

% Страницы
\usepackage{geometry}
\geometry{left=25mm,right=15mm,top=20mm,bottom=20mm}

% отступ первой строки
\usepackage{indentfirst}
\setlength{\parindent}{1.25cm}

% таблицы, картинки, списки, подписи
\usepackage{graphicx}
\usepackage{caption}
\usepackage{subcaption}
\captionsetup[table]{position=top}    % подпись таблицы над ней
\captionsetup[figure]{position=bottom}% подпись рисунка под ним
\usepackage{booktabs}
\usepackage{tabularx}
\usepackage{longtable}
\usepackage{array}
\usepackage{enumitem}
\setlist[itemize]{left=0pt,label=\textbullet}
\setlist[enumerate]{left=0pt,label=\arabic*.}

% код и конфигурации
\usepackage{listings}
\lstset{basicstyle=\ttfamily\small,breaklines=true}

% гиперссылки
\usepackage{hyperref}
\hypersetup{colorlinks=true,linkcolor=black,urlcolor=blue}

\pagestyle{plain}  % только номер страницы внизу по центру

% нумерация разделов — стандартная десятичная (article по умолчанию)
% дополнительные удобства
\usepackage{tocloft} % управление оглавлением
\setlength{\cftbeforesecskip}{6pt}

% небольшие утилиты
\newcommand{\reqfield}[1]{\textbf{#1}:} % для шаблонов в приложениях

\begin{document}

% =============================================================================
% ТИТУЛЬНЫЙ ЛИСТ
% =============================================================================
% titlepage.tex — титульный лист (первая часть названия жирная)
\begin{titlepage}
  \thispagestyle{empty} % убирает номер/колонтитул на этой странице
  \centering
  \vspace*{5.8cm}
  {\Huge \textbf{\DocumentTitleMain}\par}
  {\Large \vspace{4pt} \DocumentTitleSub \par} % вторая часть — обычным шрифтом, чуть ниже
  \vspace{0.8cm}

  {\large Версия: \Version \par}
  {\large Дата: \DocDate \par}
  \vspace{1.2cm}

  % ----- Таблица с ролями -----
  \centering
  \begin{tabularx}{0.95\textwidth}{@{\hspace{6pt}}p{5.5cm} p{6cm} p{4cm}@{\hspace{6pt}}}
    \toprule
    \textbf{Роль} & \textbf{ФИО} & \textbf{Подпись / Дата} \\
    \midrule
    Разработал: & \AuthorName & \\
    \addlinespace[6pt]
    Согласовано: & \AgreedBy & \\
    \addlinespace[6pt]
    Утверждено: & \ApprovedBy & \\
    \bottomrule
  \end{tabularx}

  \vfill
  {\large г. \City, \the\year \par}
\end{titlepage}

% =============================================================================
% ЖУРНАЛ ИЗМЕНЕНИЙ
% =============================================================================
% changes.tex — журнал изменений (образец)
\section*{Лист регистрации изменений \texttt{[MUOSD.CHANGES]}}
\addcontentsline{toc}{section}{Лист регистрации изменений}

\begin{table}[h!]
\centering
\small
\begin{tabularx}{\textwidth}{
  >{\centering\arraybackslash}p{1.1cm}
  >{\centering\arraybackslash}p{2.1cm}
  p{2.3cm}
  p{2.6cm}
  p{1.6cm}
  X
  p{1.8cm}
}
\toprule
\textbf{Версия} & \textbf{Дата} & \textbf{Автор} & \textbf{Изменённый элемент} & \textbf{Х-р изменения} & \textbf{Описание изменения} & \textbf{Утвердил} \\
\midrule
1.0 & 17.10.2025 & Фамилия И.О. & Весь документ & Добавлено & Первичный выпуск документа. & Фамилия И.О. \\
\bottomrule
\end{tabularx}
\end{table}

\clearpage
\tableofcontents
\clearpage

% =============================================================================
% ОСНОВНОЕ СОДЕРЖАНИЕ
% =============================================================================

% содержательная часть (модули)
\section{Раздел \texttt{[SECTIONID]}}
\subsection{Подраздел \texttt{[SUBSECTIONID]}}
Умный текст

\subsection{Еще один подраздел \texttt{[SUBSECTIONID]}}
Еще умного текста

\section{Раздел \texttt{[SECTIONID]}}
\subsection{Подраздел \texttt{[SUBSECTIONID]}}
Умный текст

\subsection{Еще один подраздел \texttt{[SUBSECTIONID]}}
Еще умного текста

% =============================================================================
% ПРИМЕРЫ ОФОРМЛЕНИЯ
% =============================================================================

\section{Примеры оформления рисунков}

% figures.tex — шаблоны для оформления рисунков

% =============================================================================
% ПРИМЕР ВСТАВКИ ОДИНОЧНОГО РИСУНКА С ПОДПИСЬЮ
% =============================================================================
\begin{figure}[h!]
    \centering
    \includegraphics[width=0.8\textwidth]{example-image}
    \caption{Подробное описание содержания рисунка.}
    \label{fig:example}
\end{figure}

% =============================================================================
% ШАБЛОН ДЛЯ НЕСКОЛЬКИХ РИСУНКОВ (SUBFIGURES)
% =============================================================================
\begin{figure}[h!]
    \centering
    
    \begin{minipage}{0.45\textwidth}
        \centering
        \includegraphics[width=\textwidth]{example-image-a}
        \subcaption{Описание первого рисунка}
        \label{fig:example-a}
    \end{minipage}
    \hfill
    \begin{minipage}{0.45\textwidth}
        \centering
        \includegraphics[width=\textwidth]{example-image-b}
        \subcaption{Описание второго рисунка}
        \label{fig:example-b}
    \end{minipage}
    
    \caption{Общее описание группы связанных рисунков}
    \label{fig:group}
\end{figure}

\section{Примеры оформления таблиц}

% tables.tex — шаблоны для оформления таблиц

% =============================================================================
% БАЗОВАЯ ТАБЛИЦА С ЗАГОЛОВКАМИ
% =============================================================================
\begin{table}[h!]
    \centering
    \caption{Название таблицы}
    \label{tab:example}
    
    \begin{tabularx}{0.9\textwidth}{lXl}
        \toprule
        \textbf{Колонка 1} & \textbf{Колонка 2} & \textbf{Колонка 3} \\
        \midrule
        Строка 1  & Описание элемента 1                               & 100 \\
        Строка 2  & Длинное описание элемента 2, которое может 
                    занимать несколько строк                         & 200 \\
        Строка 3  & Описание элемента 3                               & 300 \\
        \bottomrule
    \end{tabularx}
\end{table}

% =============================================================================
% ТАБЛИЦА С ГРУППИРОВКОЙ СТОЛБЦОВ
% =============================================================================
\begin{table}[h!]
    \centering
    \caption{Сложная таблица с группировкой}
    \label{tab:complex}
    
    \begin{tabular}{lcccc}
        \toprule
                          & \multicolumn{2}{c}{\textbf{Группа A}} & \multicolumn{2}{c}{\textbf{Группа B}} \\
        \cmidrule(lr){2-3} \cmidrule(lr){4-5}
        \textbf{Элемент}  & Значение 1      & Значение 2      & Значение 1      & Значение 2      \\
        \midrule
        Элемент 1         & 10              & 20              & 30              & 40              \\
        Элемент 2         & 15              & 25              & 35              & 45              \\
        Элемент 3         & 12              & 22              & 32              & 42              \\
        \bottomrule
    \end{tabular}
\end{table}

\end{document}